% ============================================================
%  附录 A  同调镜像对称导引（幻灯片）
% ============================================================
\section{同调镜像对称导引（2019 年 IMSA 报告）}

\noindent\emph{底本：D. Auroux，\emph{An invitation to homological mirror symmetry}，IMSA 成立大会，迈阿密，2019 年 9 月 6 日，共 14 页。本附录按幻灯片顺序译出，页脚标注「原稿第 $N$ 页」指幻灯片的页码。原稿第 13 页有一幅 Jeff Koons 的雕塑照片（\emph{Balloon Dog}），版权所限不予复制，正文以文字说明其位置。}

\subsection*{第 1 页：标题}
\origpage{1}

{\bfseries 同调镜像对称导引}\quad Denis Auroux，哈佛大学。IMSA 成立大会，迈阿密，2019 年 9 月 6 日。部分受 NSF 与 Simons 基金会（同调镜像对称 Simons 合作项目）资助。

\subsection*{第 2 页：Jacobi theta 函数与三角形计数}
\origpage{2}

\paragraph{椭圆曲线 $E=\C/(\Z+\tau\Z)$ 上的 Jacobi theta 函数.} 双周期的全纯函数只能是常数，但可以要求\textbf{拟周期}：
\[
s(z+1)=s(z),\qquad s(z+\tau)=e^{-\pi i\tau-2\pi iz}s(z).
\]
这样的函数在相差一个常数的意义下唯一：
\[
s(z)=\vartheta(\tau;z)=\sum_{n\in\Z}\exp\bigl(\pi in^2\tau+2\pi inz\bigr).
\]
（Jacobi，1820 年代。）

\paragraph{$T^2=\R^2/\Z^2$ 中的三角形计数（按面积加权）.} 设 $L_0,L_1,L_x$ 是三条斜率分别为 $0,1,x$ 的直线。计算
\[
?=\cdots+T^{(x-1)^2/2}+T^{x^2/2}+T^{(x+1)^2/2}+\cdots
=T^{x^2/2}\sum_{n\in\Z}T^{\frac12n^2+nx}
=e^{\pi i\tau x^2}\vartheta(\tau;\tau x),
\]
其中 $T=e^{2\pi i\tau}$ 是面积权重。

\begin{figure}[htbp]
\centering
\begin{tikzpicture}[scale=0.85,>=Stealth]
  % 左：环面上的三条闭曲线
  \begin{scope}
    \draw[thick] (0,0) ellipse (1.5 and 1.0);
    \draw[thick,blue!70!black] (-0.9,0.1) ellipse (0.85 and 0.35);
    \draw[thick,red!75!black] (-0.2,0.3) .. controls (0.9,0.55) and (1.0,-0.3) .. (0.2,-0.35);
    \draw[thick,green!45!black] (-0.4,-0.5) .. controls (0.4,-0.75) and (0.8,0.1) .. (0.5,0.35);
    \node[blue!70!black] at (-1.0,0.55) {\footnotesize $L_0$};
    \node[red!75!black] at (1.0,0.6) {\footnotesize $L_1$};
    \node[green!45!black] at (0.35,-0.85) {\footnotesize $L_x$};
  \end{scope}
  % 右：正方形里的格子与三角形
  \begin{scope}[xshift=4.2cm]
    \draw[very thick,green!45!black] (-0.2,-1.0) -- (1.8,1.4);
    \foreach \x in {0,0.8,1.6}{
      \draw[thick] (\x,-1.1) -- (\x,1.5);
    }
    \foreach \y in {-0.6,0.2,1.0}{
      \draw[thick] (-0.2,\y) -- (1.8,\y);
    }
    \draw[very thick,red!75!black] (0,0.2) -- (1.6,0.2);
    \draw[thick,blue!70!black] (0,-0.6) -- (1.6,-0.6);
    \fill (0.8,0.2) circle (1.6pt);
    \fill (0.8,0.2) circle (1.6pt);
    \node[above] at (0.8,0.25) {\footnotesize $e_0$};
    \node[below] at (0.8,-0.62) {\footnotesize $e_1$};
    \node[below left] at (0.8,0.15) {\footnotesize $s$};
  \end{scope}
\end{tikzpicture}
\caption{左：环面上的三条直线 $L_0,L_1,L_x$；右：在基本域中的三角形，其面积给出 $T$ 的幂次。（原稿第 2 页）}
\label{fig:theta}
\end{figure}

\subsection*{第 3 页：同调镜像对称（Kontsevich，1994）}
\origpage{3}

\begin{center}
\begin{tabular}{m{0.44\textwidth} m{0.44\textwidth}}
\toprule
代数（或解析）几何 & 辛几何 \\
\midrule
凝聚层（例如 $\OO_V$、向量丛 $E\to V$、摩天层 $\OO_{p\in V}$……） & $(X,\omega)$ 局部同构于 $(\R^{2n},\sum dx_i\wedge dy_i)$，拉格朗日子流形 $L$（$\dim L=n$，$\omega|_L=0$）\\
态射（及其扩张）：$H^*\hom(E,F)=\Ext^*(E,F)$ & 交点（模 Hamilton 同痕）给出 Floer 上同调 $CF^*(L,L')=K^{|L\cap L'|}$\\
导出范畴：复形 $0\to\cdots\to E^i\xrightarrow{d^i}E^{i+1}\to\cdots\to0$ 模同伦 & 微分 $\partial p=T^a q$，其中 $a=\int\omega$（并张量上局部系统数）\\
例如：函数、交点、上同调…… & 乘法 $CF(L',L'')\otimes CF(L,L')\to CF(L,L'')$，$p'\cdot p=T^a q$ \\
\bottomrule
\end{tabular}
\end{center}

两栏之间是镜像：
\[
\db\Coh(V)\ \simeq\ D^\pi\FF(X,\omega).
\]

\subsection*{第 4 页：例——椭圆曲线（Polishchuk--Zaslow）}
\origpage{4}

取
\[
E=\C/(\Z+\tau\Z),\qquad
L=\C^2\big/(z,v)\sim(z+1,v)\sim\bigl(z+\tau,e^{-\pi i\tau-2\pi iz}v\bigr),
\]
则 $\dim H^0(E,L)=1$，其截面为
\[
s(z)=\vartheta(\tau;z)=\sum_{n\in\Z}\exp\bigl(\pi in^2\tau+2\pi inz\bigr).
\]
在辛侧取 $X=T^2=\R^2/\Z^2$，$L_0,L_1,L_x$ 同上，$e_0,e_1$ 为基。此时
\[
e_1\cdot s=?\ e_0,
\]
而 $e_0$ 对应取值映射 $\OO\to\OO_x$。整个复合 $L_0\xrightarrow{s}L_1\xrightarrow{e_1}L_x$ 的计数给出与第 2 页相同的 theta 级数。

\subsection*{第 5 页：走向一般框架}
\origpage{5}

\begin{enumerate}[label=\arabic*)]
  \item 射影 Calabi--Yau 簇（$c_1=0$）：$T^2$（Polishchuk--Zaslow）、$T^{2n}$（Kontsevich--Soibelman、Fukaya、Abouzaid--Smith）、$K3$ 曲面（Seidel、Sheridan--Smith）、$X_{d=n+2}\subset\CP^{n+1}$（Sheridan）等。
  \item Fano 情形：$\CP^n$、del Pezzo 曲面、环簇等，用 LG 模型（Kontsevich、Seidel、Auroux--Katzarkov--Orlov、Abouzaid、FOOO 等）。
  \item 一般型情形、仿射簇等：紧或非紧的黎曼面（Seidel、Efimov；Abouzaid--Auroux--Efimov--Katzarkov--Orlov、Lee 等）；$\left(\C^*\right)^n$ 或环簇中的超曲面（Gammage--Shende、Abouzaid--Auroux 等）；以及其他。
\end{enumerate}
本报告的目标是给出这一纲领的一个印象，从黎曼面出发讲清 HMS 本身，不涉及 SYZ、骨骨架（skeleton）、族 Floer 理论等发展。

\subsection*{第 6 页：例 1，$F^c\simeq\db$（经典情形）}
\origpage{6}

取 $X=\R\times S^1$，$\omega=dr\wedge d\theta$，$L_r=\{r\}\times S^1$（加局部系统 $\xi$）。则
\[
HF^*(L_r,L_r)\simeq H^*(S^1,K),\qquad HF^*(L_r,L_r')=0\ (r\ne r').
\]
带局部系统的 $L_r$ 的模空间
\[
M_{\mathrm{pt}}=\bigl\{(L_r,\xi)\in\FF(X)\bigr\}/\sim
\]
有自然的解析结构，坐标为 $z(L_r,\xi)=T^r\mathrm{hol}(\xi)\in K^*$。对任意 $L'$，$CF\bigl((L_r,\xi),L'\bigr)$ 解析地依赖于 $z$。于是
\[
(L_r,\xi)\in\FF(X,\omega)\ \longleftrightarrow\ \OO_z\in\db\bigl(X^\vee=K^*\bigr).
\]
这就是 SYZ 图景：$X$ 为 Calabi--Yau，$\pi\colon X\to B$ 为拉格朗日环面纤维化，则
\[
X^\vee=\bigl\{\OO_p,\ p\in X^\vee\bigr\}=\bigl\{(L_b=\pi^{-1}(b),\xi)\in\FF(X)\bigr\}/\sim .
\]

\subsection*{第 7 页：例 1，$F^{\mathrm{wr}}\simeq\db$（Abouzaid--Seidel）}
% [Tutor Note: 易错] 第 6 页与第 7 页的差别只在范畴的选择：紧拉格朗日里 $L_r$ 与 $L_{r'}$ 不相交时 $\hom=0$，所以只能靠 $L_r$ 里的局部系统 $\xi$ 记录 $z$；换成 wrapped 之后，$L_0$ 可以沿 Reeb 方向卷绕，于是单个对象 $L_0$ 就承载了整个 $K[x^{\pm1}]$。两条路线给出同一个 $\db\Coh(X^\vee)$，但后者的生成性是一句话的事。来源：附录第 6、7、8 页。
\origpage{7}

考虑 wrapped Fukaya 范畴。取 $X=\R\times S^1\supset L_0=\R\times\{0\}$，后者是非紧拉格朗日。取 Hamilton 扰动 $H=\tfrac12r^2$，则
\[
\phi^1_H(r,\theta)=(r,\theta+r),
\]
交点出现在 $X$ 的内部以及边界处的 Reeb 流上。定义
\[
CW^*(L_0,L_0):=CF^*\bigl(\phi^1_H(L_0),L_0\bigr)=\bigoplus_{i\in\Z}Kx_i .
\]
乘法由三角形计数给出：$\tilde q\in\phi^2_H(L)\cap L$ 与 $q\in\phi^1_H(L)\cap L$ 通过 $r\mapsto2r$ 对应。于是
\[
x_k\cdot x_l=x_{k+l},\qquad \End(L_0)\simeq K\bigl[x^{\pm1}\bigr],
\]
其中 $x^k\rightsquigarrow x^k$。

\subsection*{第 8 页：例 1，$F^{\mathrm{wr}}\simeq\db(K^*)$}
\origpage{8}

承上，$\End(L_0)\simeq K[x^{\pm1}]\simeq\End(\OO_{X^\vee})$。并且 $L_0$ 生成 $F^{\mathrm{wr}}(X)$。由 Yoneda 引理，$L\mapsto\Hom(L_0,L)$ 给出嵌入
\[
F^{\mathrm{wr}}(X)\hookrightarrow \End(L_0)\text{-}\mathrm{mod}.
\]
例如
\[
(L_r,\xi)\ \longmapsto\ HF\bigl(L_0,(L_r,\xi)\bigr)\simeq K\bigl[x^{\pm1}\bigr]/(x-z),
\qquad z=T^r\mathrm{hol}(\xi).
\]
\begin{theo}
$F^{\mathrm{wr}}(X)\simeq K[x^{\pm1}]\text{-}\mathrm{mod}\simeq\db\Coh(X^\vee)$。
\end{theo}

\subsection*{第 9 页：例 2，$F^{\mathrm{wr}}$（AAEKO）}
\origpage{9}

取
\[
X=S^2\setminus\{-1,0,\infty\}=\C^*\setminus\{-1\},\qquad L_0=\R_+ .
\]
则 $CW(L_0,L_0)=\bigoplus_{i\in\Z}Kx_i$，乘法为
\[
x_j\cdot x_i=\begin{cases}x_{i+j}, & ij\ge0,\\ 0, & ij<0,\end{cases}
\qquad\End(L_0)\simeq K[x,y]/(xy=0).
\]
于是自然要问：镜像是否为
\[
X^\vee=\Spec K[x,y]/(xy=0)=\{xy=0\}\subset\A^2,
\]
以及嵌入 $F^{\mathrm{wr}}(X)\hookrightarrow\End(L_0)\text{-}\mathrm{mod}$ 是否成立？

\subsection*{第 10 页：例 2，$F^{\mathrm{wr}}\simeq\db(\{xy=0\})$}
\origpage{10}

取 $X=\C^*\setminus\{-1\}$，三条射线
\[
L_0=(0,\infty),\qquad L_1=(-1,0),\qquad L_2=(-\infty,-1)
\]
生成 $F^{\mathrm{wr}}(X)$。它们的自同态环分别是
\[
K[x,y]/(xy),\qquad K[y,z]/(yz),\qquad K[x,z]/(xz),
\]
彼此之间有乘法关系 $u_y\cdot v_y=y$、$v_y\cdot u_y=y$，并满足正合三角
\[
L_2\xrightarrow{u_x}L_0\xrightarrow{u_y}L_1\xrightarrow{u_z}L_2[1],
\qquad
L_1\xrightarrow{v_y}L_0\xrightarrow{v_x}L_2\xrightarrow{v_z}L_1[1].
\]

\subsection*{第 11 页：例 2（续）}
\origpage{11}

B 侧取 $X^\vee=\{xy=0\}=A\cup B\subset\A^2$。于是
\[
\OO\ \leftrightarrow\ K[x,y]/(xy)=:R,\qquad
\OO_A\ \leftrightarrow\ K[y,z]/(yz),\qquad
\OO_B\ \leftrightarrow\ K[x,z]/(xz),
\]
并有 $\Hom(\OO_A,\OO_A)\simeq K[y]$、$\Ext^2_k(\OO_A,\OO_A)\ni z^k$。对应的正合三角为
\[
\OO_B\xrightarrow{x}\OO\xrightarrow{1}\OO_A\xrightarrow{u}\OO_B[1],
\qquad
\OO_A\xrightarrow{y}\OO\xrightarrow{1}\OO_B\xrightarrow{v}\OO_A[1].
\]
\begin{theo}[AAEKO]
$F^{\mathrm{wr}}(X)\simeq\db\Coh(X^\vee)$。
\end{theo}

\subsection*{第 12 页：例 2，$F^{\mathrm{wr}}\simeq\db_{\mathrm{sing}}(\C^3,-xyz)$}
\origpage{12}

取 $X=\PP^1\setminus\{-1,0,\infty\}$，$X^\vee=\{xy=0\}$。此时 $F^{\mathrm{wr}}(X)\simeq\db\Coh(\{xy=0\})$ 关于 $x,y,z$ 并不对称，无法直接推广到更高亏格，需要粘合。

\paragraph{稳定化.} 把 $X$ 写成 $X\simeq\{x+y+1=0\}\subset\left(\C^*\right)^2$。取
\[
\Bigl(\mathbb X=\mathrm{Bl}\bigl(\left(\C^*\right)^2\times\C,\ X\times0\bigr),\ W=p_{\C}\Bigr)
\ \longleftrightarrow\
\bigl(X^\vee=\C^3,\ W^\vee=-xyz\bigr).
\]
\begin{theo}[AAEKO]
\[
F^{\mathrm{wr}}(X)\simeq\db_{\mathrm{sing}}(X^\vee,W^\vee):=\db\Coh(\{xyz=0\})/\Perf
\]
（Orlov 的商），并且 $(L_0,L_1,L_2)\leftrightarrow\bigl([\OO_{\{z=0\}}],[\OO_{\{x=0\}}],[\OO_{\{y=0\}}]\bigr)$。
\end{theo}
这一结果可以推广到所有黎曼面（AAEKO、Seidel、Efimov、H. Lee）。镜像 $(X^\vee,W^\vee)$ 的维数为 $3$（Hori--Vafa、AAK）。

\subsection*{第 13 页：$(X^\vee,W^\vee)$ 的几何}
\origpage{13}

对仿射平面曲线 $\Sigma=\{f(x,y)=0\}\subset\left(\C^*\right)^2$，其镜像为：$X^\vee$ 是由 $f$ 的热带化确定的环 Calabi--Yau 三维簇，$W^\vee\in\OO(X^\vee)$，而
\[
Z:=\{W^\vee=0\}=\cup\ \text{环层},\qquad
\mathrm{sing}(Z)=\crit(W^\vee)=\cup\ \text{一维环层}=\PP^1\ \text{与}\ \A^1\ \text{之并}.
\]
镜像分解为
\[
\Sigma=\cup\ \{\text{裤子}\}\ \longleftrightarrow\ (X^\vee,W^\vee)=\cup\ (\C^3,-xyz).
\]
原幻灯片此处配有一幅 Jeff Koons 的雕塑 \emph{Balloon Dog} 照片（摄影：Librado Romero，\emph{The New York Times}），用以图解"由若干块粘成的气球狗"与"由若干 $(\C^3,-xyz)$ 粘成的镜像"之间的类比。版权所限，本译本不复制该照片。

\begin{theo}[Heather Lee]
\[
F^{\mathrm{wr}}(\Sigma)\simeq\varinjlim\Bigl(F^{\mathrm{wr}}(\text{各块})\Rightarrow F^{\mathrm{wr}}(\text{整体})\Bigr)
\simeq\db_{\mathrm{sing}}(X^\vee,W^\vee)=\db(Z)/\Perf .
\]
\end{theo}
相关工作：Bocklandt、Gammage--Shende、Lekili--Polishchuk 等。

\subsection*{第 14 页：逆命题}
\origpage{14}

\begin{theo}[Abouzaid--A.]
逆命题也成立：
\[
\FF(X^\vee,W^\vee)\ \simeq\ \db\Coh(\Sigma).
\]
\end{theo}
A.--Efimov--Katzarkov 正在进行的工作把左端改写成 $\crit(W^\vee)=Z$ 的一维层上的对象；曲线在阿贝尔曲面中的情形见 C. Cannizzo 的学位论文。Abouzaid--A. 的结果对 $X$ 为 $\left(\C^*\right)^n$ 中的超曲面或完全交也成立。

\begin{remark}
附录与正文的关系可以这样看：正文第 1--3 讲给出 $H\subset\left(\C^*\right)^n$ 的镜像与 $F(Y,W)$ 的构造，本附录则把同一套机器放到 $\Sigma\subset\left(\C^*\right)^2$ 的最简单情形，用三条射线、柱面与椭圆曲线把它算穿。第 12 页的 $(\C^3,-xyz)$ 正是正文第 1 页例~\ref{ex:pants} 的 $\bigl(\C^3,-z_1z_2z_3\bigr)$；第 14 页的 $\db\Coh(\Sigma)$ 正是正文第 16 页定理~\ref{thm:aek} 的右端。
% [Tutor Note: 点拨] 正文与附录的记号有一处差异：附录用 $X,X^\vee$ 记镜像对，正文用 $V,Y$。两处 $(\C^3,-xyz)$ 是同一个对象。来源：原书第 1 页与附录第 12 页。
\end{remark}
